\subsection{An adaptive lower tail for batch collisions}
\label{app:lower-proof-collision-tail}

The conditional Laplace bound and the edge-counting identity now force many collisions
with constant probability under the joint instance--midpoint law. This is
the input needed to accumulate edge-level discrepancy nonlinearly; no
independence between commit rounds is assumed. Recall that
\(\TotalCollisions=\sum_t\CollisionCount{t}\) counts the unresolved edges
whose two endpoints are committed in the same active batch, as defined in
\eqref{app:eq:unresolved-batch-and-collisions}.

\paragraph{Intuition.}
If order $N$ unresolved vertices were spread evenly across $\RoundBudget$
batches while the unresolved pool had size order $N$, the collision scale
would be
\[
 \RoundBudget\frac{(N/\RoundBudget)^2}{N}
 =\frac{N}{\RoundBudget}.
\]
The proof does not assume even batches or a fixed pool: the edge-counting identity
forces enough total active batch mass, Cauchy--Schwarz gives the same
$N/\RoundBudget$ scale, and the conditional Laplace bound turns that
pathwise mass into a constant-probability collision event.

\begin{proposition}[Adaptive collision lower tail]
\label{app:prop:collision-lower-tail}
Fix a decision-rule seed under the joint law
\eqref{app:eq:joint-instance-midpoint-law}. If
\(\QueryBudget{}\le \HardBankSize/8\) and \(1\le \RoundBudget{}\le N/8\), then
\begin{equation}
 \mathbb P\left(
   \TotalCollisions\ge\frac{N}{8192\RoundBudget{}}
 \right)
 \ge
 \frac12-\exp\left\{-\frac{N}{8192\RoundBudget{}}\right\}.
 \label{app:eq:collision-lower-tail}
\end{equation}
In particular, if \(\RoundBudget{}\le N/16384\), then
\begin{equation}
 \mathbb P\left(
   \TotalCollisions\ge\frac{N}{8192\RoundBudget{}}
 \right)\ge\frac14.
 \label{app:eq:collision-quarter-probability}
\end{equation}
\end{proposition}

\begin{proof}

\emph{1. Fix the pre-exposure information and conditional Laplace bound.}
For an actual commit round \(t\), let \(B_t\) be the batch selected by the
seed-fixed decision rule. Recall the existing notation:
\[
 \ActiveBatch{t}=B_t\cap\UnresolvedVertices{t},
 \qquad
 \UnresolvedVertexCount{t}=|\UnresolvedVertices{t}|,
 \qquad
 \ActiveBatchSize{t}=|\ActiveBatch{t}|,
\]
\[
 \CollisionCount{t}
 =|\{e\in\OpUnresolvedMatching{s_t}:
                    e\subseteq\ActiveBatch{t}\}|,
 \qquad
 \TotalCollisions=\sum_t\CollisionCount{t}.
\]
Thus \(B_t\) is the full commit batch, whereas \(\ActiveBatch{t}\)
contains only its currently unresolved vertices.

After termination, append empty rounds until the total number of indexed
rounds is \(\RoundBudget{}\). In every appended round set
\(B_t=\ActiveBatch{t}=\UnresolvedVertices{t}=\varnothing\), so
\(\UnresolvedVertexCount{t}=\ActiveBatchSize{t}=\CollisionCount{t}=0\).
These rounds add no collision, and hence
\(\TotalCollisions=\sum_{t=1}^{\RoundBudget{}}\CollisionCount{t}\).
Introduce only the following additional abbreviation:
\[
 A_t:=
 \begin{cases}
  \PositivePart{\ActiveBatchSize{t}-1}^2/\UnresolvedVertexCount{t},
      &\UnresolvedVertexCount{t}>0,\\
  0,&\UnresolvedVertexCount{t}=0.
 \end{cases}
\]
The pre-commit analytic field \(\AnalyticField{t}\) determines both the
selected batch \(B_t\) and the unresolved set \(\UnresolvedVertices{t}\),
hence also \(\ActiveBatch{t}=B_t\cap\UnresolvedVertices{t}\).
Selecting \(B_t\) adds no observation: the matching edges incident to
\(\ActiveBatch{t}\) and the midpoint values \(X_{B_t}\) have not yet been
exposed. Thus \(\UnresolvedVertexCount{t}\), \(\ActiveBatchSize{t}\),
\(A_t\), and all quantities from earlier rounds are
\(\AnalyticField{t}\)-measurable. For appended empty rounds, use the terminal
field.

By \cref{app:lem:unresolved-invariant}, conditional on \(\AnalyticField{t}\),
the matching on \(\UnresolvedVertices{t}\) is uniform. For
\(\UnresolvedVertexCount{t}\ge2\), apply
\cref{app:lem:matching-collision-laplace} with
\(U=\UnresolvedVertices{t}\), \(S=\ActiveBatch{t}\),
\(n=\UnresolvedVertexCount{t}\), \(k=\ActiveBatchSize{t}\), and Laplace
parameter one. Since \(1-e^{-1}\ge1/2\), this gives
\begin{equation}
 \mathbb E\!\left[e^{-\CollisionCount{t}}\mid\AnalyticField{t}\right]
 \le
 \exp\left\{-\frac{1-e^{-1}}{32}A_t\right\}
 \le e^{-A_t/64}.
 \label{app:eq:conditional-collision-laplace}
\end{equation}
The unresolved set has even cardinality, so
\(\UnresolvedVertexCount{t}=1\) cannot occur. If
\(\UnresolvedVertexCount{t}=0\), then
\(\ActiveBatchSize{t}=\CollisionCount{t}=A_t=0\), and the same display holds
with equality.


\emph{2. Iterate the exponential expectation bound.}
Let \(r=0,\ldots,\RoundBudget{}\) count completed commit rounds, not
individual operations.  Define
\begin{equation}
 \CollisionExponential{r}:=
 \exp\left\{
   -\sum_{t=1}^{r}\CollisionCount{t}
   +\frac1{64}\sum_{t=1}^{r}A_t
 \right\},
 \qquad \CollisionExponential{0}:=1.
 \label{app:eq:collision-supermartingale}
\end{equation}
Both \(\CollisionExponential{r-1}\) and \(A_r\) are measurable with respect
to \(\AnalyticField{r}\), the field immediately before commit \(r\)
after any intervening counterfactual operations. Hence
\eqref{app:eq:conditional-collision-laplace} and the tower property give
\[
 \mathbb E \CollisionExponential{r}
 =
 \mathbb E\!\left[
   \CollisionExponential{r-1}e^{A_r/64}
   \mathbb E(e^{-\CollisionCount{r}}\mid\AnalyticField{r})
 \right]
 \le \mathbb E \CollisionExponential{r-1}.
\]
Iteration yields
\begin{equation}
 \mathbb E \CollisionExponential{\RoundBudget}\le1.
 \label{app:eq:terminal-collision-supermartingale}
\end{equation}
This is the only concentration step; adaptivity is absorbed into the
successive conditional expectations.


\emph{3. Turn edge counting into a pathwise Laplace budget.}
By \cref{app:lem:query-disclosure} and
\(\QueryBudget{}\le\HardBankSize/8\),
\[
 \mathbb E\GenieDeletions
 \le\frac{N\QueryBudget{}}{\HardBankSize}
 \le\frac{N(\HardBankSize/8)}{\HardBankSize}
 =\frac N8.
\]
Because \(\GenieDeletions\ge0\), Markov's inequality gives
\[
 \mathbb P(\GenieDeletions>N/4)
 \le\frac{\mathbb E\GenieDeletions}{N/4}
 \le\frac{N/8}{N/4}=\frac12.
\]
Taking the complement,
\(\mathbb P(\GenieDeletions\le N/4)=1-\mathbb P(\GenieDeletions>N/4)\),
therefore yields
\begin{equation}
 \mathbb P(\GenieDeletions\le N/4)\ge\frac12.
 \label{app:eq:few-genie-deletions-event}
\end{equation}
On this event, the pathwise counting identity
\eqref{app:eq:active-excess-ledger} in \cref{app:lem:pathwise-edge-counting}, the bounds
\(\ActiveRoundCount\le \RoundBudget{}\le N/8\), and
\(\TotalCollisions\ge0\) imply
\begin{equation}
 \sum_{t=1}^{\RoundBudget{}}\PositivePart{\ActiveBatchSize{t}-1}
 \ge \frac N2-\GenieDeletions-\RoundBudget{}
 \ge\frac N8.
 \label{app:eq:pathwise-active-excess}
\end{equation}
Since \(\UnresolvedVertexCount{t}\le N\), the definition of \(A_t\)
implies \(A_t\ge\PositivePart{\ActiveBatchSize{t}-1}^2/N\).
This also holds when \(\UnresolvedVertexCount{t}=0\), because then both
sides are zero. Cauchy--Schwarz over the \(\RoundBudget{}\) indexed rounds
(including appended empty rounds) now gives
\begin{align}
 \sum_{t=1}^{\RoundBudget{}}A_t
 &\ge\frac1N\sum_{t=1}^{\RoundBudget{}}\PositivePart{\ActiveBatchSize{t}-1}^2\notag\\
 &\ge\frac1{N\RoundBudget{}}
       \left(\sum_{t=1}^{\RoundBudget{}}\PositivePart{\ActiveBatchSize{t}-1}\right)^2
 \ge\frac{N}{64\RoundBudget{}}.
 \label{app:eq:pathwise-laplace-budget}
\end{align}


\emph{4. Subtract the small-collision event.}
Set \(x:=N/(8192\RoundBudget{})\). On the intersection
\(\{\GenieDeletions\le N/4\}\cap\{\TotalCollisions<x\}\),
\eqref{app:eq:pathwise-laplace-budget} yields
\[
 \log \CollisionExponential{\RoundBudget}
 =-\TotalCollisions+\frac1{64}\sum_{t=1}^{\RoundBudget{}}A_t
 >-x+\frac{N}{4096\RoundBudget{}}
 =x.
\]
Using \eqref{app:eq:terminal-collision-supermartingale} and Markov's
inequality,
\begin{align*}
 \mathbb P(\GenieDeletions\le N/4,\ \TotalCollisions<x)
 &\le \mathbb P(\CollisionExponential{\RoundBudget}>e^x)\\
 &\le\frac{\mathbb E\CollisionExponential{\RoundBudget}}{e^x}
 \le\frac1{e^x}=e^{-x}.
\end{align*}
Combining with \eqref{app:eq:few-genie-deletions-event},
\begin{align*}
 \mathbb P(\TotalCollisions\ge x)
 &\ge \mathbb P(\GenieDeletions\le N/4,\ \TotalCollisions\ge x)\\
 &=\mathbb P(\GenieDeletions\le N/4)
   -\mathbb P(\GenieDeletions\le N/4,\ \TotalCollisions<x)\\
 &\ge \frac12-e^{-x},
\end{align*}
which proves \eqref{app:eq:collision-lower-tail}.
If \(\RoundBudget{}\le N/16384\), then \(x\ge2\)
and \(1/2-e^{-2}>1/4\), proving
\eqref{app:eq:collision-quarter-probability}.
\end{proof}
